
\subsubsection{2.1 Gradient-Based Adaptive Rank Methods}

The dominant approach to per-layer rank selection uses gradient signals during training to estimate layer importance.

\textbf{AdaLoRA} \citep{Zhang2023AdaLoRA} parameterizes weight updates in singular value decomposition form and prunes singular values based on importance scores derived from gradient magnitudes. Rank allocation is dynamic, changing during training. The singular values being pruned belong to $\Delta W$ (the learned update), not $W_0$ (the pretrained weights) --- a fundamental distinction from the approach investigated here.

\textbf{DyLoRA} \citep{Valipour2022} trains LoRA adapters simultaneously across a range of ranks, enabling rank selection post-training by truncating the adapter. This eliminates grid search but still requires full training before the optimal rank is known.

\textbf{LAARA} \citep{Tripathi2026LAARA} provides a formal theoretical foundation for per-layer allocation using Fisher information and proves that uniform rank is suboptimal. Its implementation uses gradient warmup to estimate layer importance. LAARA is the closest theoretical peer to the structural approach investigated here.

\textbf{IGU-LoRA} \citep{Jiang2026IGULoRA} identifies a gradient bias in AdaLoRA's importance scores and proposes integrated gradient corrections.

\subsubsection{2.2 Calibration-Based Methods}

\textbf{IFCLoRA} \citep{Zhang2026IFCLoRA} computes per-layer importance from Information Flow Centrality scores derived from a calibration forward pass. IFCLoRA demonstrates that pre-fine-tuning signals can predict rank need. The key distinction from the structural approach: IFCLoRA requires calibration data and a forward pass; the present method requires only the pretrained weight matrices.

\subsubsection{2.3 Structural Methods Using $W_0$}

\textbf{PiSSA} \citep{Meng2024PiSSA} initializes LoRA matrices from the principal singular vectors of $W_0$, achieving faster convergence. Rank is still set uniformly; the SVD of $W_0$ informs initialization, not rank selection.

\textbf{LoRA-XS} \citep{Balazy2024LoRAXS} freezes the full $W_0$ SVD as a structured matrix and trains only a small $r \times r$ core adapter. Like PiSSA, it uses $W_0$ structure for adapter design, not rank selection.

Both works establish that $W_0$'s singular structure contains adaptation-relevant information but do not test whether it predicts optimal per-layer rank.

\subsubsection{2.4 Intrinsic Dimensionality and Effective Rank}

\textbf{Aghajanyan et al.} [2021] showed that the intrinsic dimensionality of fine-tuning loss landscapes varies by layer, with $d_{90}$ differing substantially between attention and FFN layers. This observation motivates the hypothesis that $W_0$ structure encodes adaptation complexity.

\textbf{Roy and Vetterli} [2007] define effective rank $\text{erank}(A) = \exp(H(\sigma/\|\sigma\|_1))$, bounded in $[1, \text{rank}(A)]$, scale-invariant, and providing wider dynamic range than stable rank or spectral norm ratios.

\subsubsection{2.5 Position}

No published work has tested whether $\text{erank}(W_0)$ --- a scalar computed from pretrained weights before any fine-tuning --- correlates significantly with per-layer PARA oracle ranks. The present work investigates this question and reports preliminary results for BERT-base-uncased.

